\section{Discussion}

Our experiments confirm that spurious features dominate early in ERM training (H-E1) but falsify the gradient competition mechanism (H-M1). We discuss the implications of this falsification, propose an alternative mechanistic explanation, and acknowledge limitations.

\subsection{Key Finding: Inverted Gradient Dynamics}

The gradient competition hypothesis assumes that spurious features dominate because they ``win'' the competition for gradient signal. Our measurements reveal the opposite: minority samples---those requiring invariant feature learning---produce 5.7$\times$ higher gradient norms than spurious-aligned samples.

This inversion makes sense when we reconsider what gradient magnitude indicates. High gradients signal \emph{ongoing learning}---samples that remain difficult continue generating parameter updates. Low gradients indicate \emph{completed learning}---samples that the model has already solved contribute little to further updates.

\subsection{Alternative Mechanism: Convergence Speed}

We propose that spurious dominance arises from \emph{convergence speed} rather than gradient magnitude:

\begin{enumerate}
    \item \textbf{Simpler patterns occupy flatter loss landscape regions} that are easier to navigate
    \item \textbf{Networks converge to spurious solutions faster} than to invariant solutions
    \item \textbf{Early convergence locks in spurious representations} before core features develop
    \item \textbf{High minority gradients cannot overcome established spurious features} because the network's effective learning rate on spurious-aligned samples is already near zero
\end{enumerate}

This reframing connects to Sharpness-Aware Minimization (SAM) \citep{Foret2021Sharpness}. If spurious solutions occupy sharp minima, SAM's sharpness penalty would naturally disfavor spurious features.

\subsection{Implications for Robustness Methods}

\textbf{Methods targeting gradient rebalancing may be misguided.} Our results show hard samples already dominate gradient flow---the problem is not gradient magnitude but convergence timing.

\textbf{Loss landscape interventions may be more effective.} Methods that target loss landscape geometry (SAM, weight averaging) rather than sample-level gradients may address the actual mechanism of spurious dominance.

\subsection{Limitations}

\textbf{Single dataset.} We test only on Waterbirds. Spurious correlation dynamics may differ on CelebA or CMNIST.

\textbf{Synthetic region masks.} Our GradCAM analysis uses spatial heuristics rather than ground-truth segmentation masks.

\textbf{Layer4 gradients only.} Earlier layers may show different dynamics.

\textbf{ResNet-50 architecture.} Vision Transformers may exhibit different spurious learning dynamics.

\subsection{Broader Impact}

\textbf{Positive impacts.} Understanding the mechanism of spurious feature learning helps develop more effective robustness interventions.

\textbf{Potential misuse.} We do not anticipate direct misuse of this research.
